\chapter{Post-review Mechanism Checks}
\label{app:mechanism}

This appendix reports an explanatory follow-up performed after the
confirmatory protocol had been frozen and run. It does not replace or modify
the 3,111 confirmatory configurations. Its purpose is to test why the two
references behave differently and whether the complete variance sensitivity
is responsible for the observed Expanded Welch adjustment.

The follow-up selected 809 distinct configurations from the saved main,
baseline, sample-imbalance and convergence families. Each configuration used
20,000 new table pairs. The 169 null configurations also used a separate
20,000-pair pilot to estimate finite-sample moments without reusing the
evaluation draws. In total, the follow-up generated 19.56 million table
pairs. The populations were read unchanged from the confirmatory population
file; the follow-up therefore changes the diagnostic calculation, not the
population construction.

\section{Local explanation near independence}

The approximation in Section~\ref{sec:near-independence-mechanism} can be
seen through a local Gaussian representation. After the row and column
effects are removed, let \(u\in\mathbb R^d\) denote the population
interaction coordinates, where \(d=(r-1)(c-1)\). In fixed dimension, with
positive cell probabilities and sufficiently large expected counts, the
standardised empirical interaction is approximately
\begin{equation}
  \sqrt n\,\widehat u \approx \sqrt n\,u+Z,
  \qquad Z\sim N_d(0,I_d).
\end{equation}
Locally, \(\lVert u\rVert^2\approx 2I\approx V\). Hence the quadratic
quantity
\begin{equation}
  R=\lVert\sqrt n\,u+Z\rVert^2
\end{equation}
is approximately noncentral chi-squared with \(d\) degrees of freedom and
noncentrality \(nV\), while \(2n\widehat I\approx R\) and
\(\widehat V\approx R/n\). The standard moments
\(\E R=d+nV\) and \(\Var(R)=2d+4nV\) then give
\begin{align}
  \E\widehat V&\approx V+d/n,\\
  \Var(\widehat I)&\approx V/n+d/(2n^2),\\
  \Var(\widehat V)&\approx 4V/n+2d/n^2.
\end{align}
These equations separate two effects. The linear approximation contributes
the terms containing \(V\); random quadratic interaction contributes the
terms containing \(d\). When \(nV\) is not large relative to \(d\), both the
estimated standard error and the component degree of freedom can therefore
depart materially from their first-order population approximations. This is
a local, fixed-dimension explanation, not an exact formula for sparse tables.

\section{Independent standard-deviation diagnostic}

For each selected null, an independent pilot estimates the actual
finite-sample standard deviation of the corrected MI difference. The
evaluation numerator is then divided by this single pilot value. This
``MC SD'' is unavailable from one real data set because it uses repeated
samples from the known simulation populations. It is used only to diagnose
the approximation. It simultaneously removes random denominator error and
numerator--denominator dependence, so it cannot attribute a discrepancy to
mean denominator bias alone.

\begin{table}[htbp]
\centering\small
\caption{Denominator diagnostic for the additive first-block $3\times3$ different-skew null, $I(P)=I(Q)=0.02$. $R$ is mean estimated squared SE divided by independently estimated sampling variance. MC SD uses that independent finite-sample SD; first-order SD uses the population influence variance. Each pilot and evaluation uses 20,000 independent replicates.}
\label{tab:denominator-check}
\begin{tabular}{rrrrr}
\toprule
$n$ & $R$ & Wald & MC SD & First-order SD\\
\midrule
100 & 1.17 & 0.0166 & 0.0551 & 0.0976 \\
1000 & 1.06 & 0.0384 & 0.0509 & 0.0532 \\
\bottomrule
\end{tabular}
\end{table}
\begin{table}[htbp]
\centering\small
\caption{Variance and component-df diagnostics for the same populations. The population first-order df is $2nV^2/\tau^2$; plug-in is its median over evaluation tables; MC is $2\overline{\widehat V}^{2}/s^2_{\widehat V}$ from the independent pilot.}
\label{tab:component-check}
\begin{tabular}{rlrrrr}
\toprule
$n$ & Population & $\E\widehat V/V$ & First order & Plug-in & MC\\
\midrule
100 & P & 1.75 & 1.73 & 3.14 & 4.84 \\
100 & Q & 1.63 & 1.51 & 2.43 & 3.77 \\
1000 & P & 1.09 & 17.26 & 18.82 & 20.52 \\
1000 & Q & 1.08 & 15.08 & 16.52 & 17.56 \\
\bottomrule
\end{tabular}
\end{table}


Table~\ref{tab:denominator-check} shows that the estimated squared standard
error is 17\% too large relative to the empirical finite-sample variance at
\(n=100\). Replacing it with the independent MC SD moves rejection from
0.0166 to 0.0551. However, the population first-order SD gives 0.0976,
because the finite-sample numerator variance is itself 42\% larger than its
first-order value. At \(n=1000\), the corresponding approximations are much
closer. Thus denominator estimation is an important mechanism in this
example, but the evidence does not support treating it as the only one.

Table~\ref{tab:component-check} makes the degree-of-freedom distinction
explicit. At \(n=100\), the population first-order component values are 1.73
and 1.51, whereas moment matching the simulated \(\widehat V\) gives 4.84
and 3.77. The variance approximation
\(\Var(\widehat V)\approx\tau^2/n\) is already within about 10\%; much of the
remaining degree-of-freedom difference comes from the upward shift in
\(\E\widehat V\). By \(n=1000\), all three component calculations are much
closer. A finite-difference derivative check validates the formula for
\(\tau^2\), but only this repeated-sampling check assesses whether its
first-order moment approximation is accurate.

\section{Ablations of the degrees of freedom}

Three simpler references were evaluated with the same corrected numerator and
plug-in standard error:
\begin{enumerate}
  \item The Hutcheson-style MI analogue sets each component degree of freedom
  to \(n\), following Equation~\eqref{eq:hutcheson-df}.
  \item Simple Welch sets each component degree of freedom to \(n-1\), as for
  ordinary sample variances.
  \item Kurtosis-only Welch freezes the pointwise scores when differentiating
  the variance. Its sensitivity is
  \(g^{(K)}_{ij}=(\ell_{ij}-I)^2-V\), giving
  \begin{equation}
    \nu^{(K)}=\frac{2nV^2}{\sum_{ij}p_{ij}\{g^{(K)}_{ij}\}^2}.
  \end{equation}
\end{enumerate}
The third ablation omits the change in the pointwise MI values and their row
and column margins. It therefore isolates what the complete derivative adds
to the reference calculation.

\begin{figure}[htbp]
\centering
\includegraphics[width=\textwidth]{mechanism_ablation.pdf}
\caption[Reference ablations in the $3\times3$ main family]{Reference
ablations in the \(3\times3\) main family. Rows are uniform, same-skew and
different-skew margins; columns use equal sample sizes 10, 100 and 1000. The
horizontal axis is the constructed MI difference from 0 to 0.02 nats and the
vertical axis is unconditional rejection from 0 to 0.15. Each point uses
20,000 new table pairs. Hollow markers indicate validity below 90\%; complete
rates, including values outside this calibration zoom, are retained in the
follow-up CSV.}
\label{fig:mechanism-ablation}
\end{figure}

Figure~\ref{fig:mechanism-ablation} shows that the Hutcheson-style and Simple
Welch references are usually close to each other and to Wald, while the
kurtosis-only correction lies between them and Expanded Welch. Their rejection
rates are identical in 494 of 809 configurations and differ by 0.00047 on
average in absolute value. The largest difference, 0.01235, occurs in the
extreme \(5\times5\), \(n_P=n_Q=5\) case, where replacing \(n_i-1=4\) by
\(n_i=5\) is proportionally substantial. The complete pointwise and margin
sensitivity is therefore what drives most of the extra tail inflation near low
MI. This confirms that the expanded derivative has a real numerical effect; it
does not establish that the effect improves calibration.

\section{Bias-correction sensitivity}

The confirmatory statistic uses the declared fixed-support correction
\(d/(2n)\). As a sensitivity check only, the follow-up replaced \(d\) by
the observed support quantity
\begin{equation}
  d_{\rm obs}=K_{XY}-K_X-K_Y+1,
\end{equation}
where the three \(K\)'s count occupied joint cells, rows and columns. In the
different-skew \(3\times3\) equal-MI imbalance family with baseline MI 0.02,
this changed Wald rejection from 0.0661 to 0.0580 when
\((n_P,n_Q)=(50,500)\), but from 0.0827 to 0.0262 when the allocation was
reversed. The sensitivity can therefore change direction asymmetry, but it
is not a validated universal correction. Its dependence on random observed
support introduces a new source of variation, so the declared-support
statistic remains the confirmatory definition.

\section{Complete main-null table}

Table~\ref{tab:all-main-null} provides the exact rejection and validity rates
for all 108 main-grid null configurations. It separates the 36 different-skew
weak-null configurations, where \(P\ne Q\) but \(I(P)=I(Q)\), from the 72
uniform and same-skew strong-null configurations, where \(P=Q\). These rows
are the source data for the summary statements in Chapter~\ref{chap:results}.

\begingroup\footnotesize
\setlength{\tabcolsep}{4pt}
\begin{longtable}{@{}llrrrrr@{}}
\caption{All 108 main null regimes in the frozen run. Each rate uses all 20,000 replicates. V denotes the fraction with a valid result. Uniform and same-skew profiles have $P=Q$; different-skew profiles have $P\ne Q$ with equal MI.}\label{tab:all-main-null}\\
\toprule Shape & Margins & $n$ & Wald & V & Expanded & V\\\midrule\endfirsthead
\toprule Shape & Margins & $n$ & Wald & V & Expanded & V\\\midrule\endhead
2x2 & different skew & 2 & 0.0000 & 0.0000 & 0.0000 & 0.0000 \\
2x2 & different skew & 5 & 0.1008 & 0.8454 & 0.0000 & 0.3423 \\
2x2 & different skew & 10 & 0.0432 & 0.9791 & 0.0120 & 0.7575 \\
2x2 & different skew & 20 & 0.0244 & 0.9983 & 0.0151 & 0.9749 \\
2x2 & different skew & 50 & 0.0080 & 1.0000 & 0.0039 & 0.9999 \\
2x2 & different skew & 100 & 0.0121 & 1.0000 & 0.0054 & 1.0000 \\
2x2 & different skew & 250 & 0.0290 & 1.0000 & 0.0158 & 1.0000 \\
2x2 & different skew & 500 & 0.0387 & 1.0000 & 0.0283 & 1.0000 \\
2x2 & different skew & 1000 & 0.0453 & 1.0000 & 0.0380 & 1.0000 \\
2x2 & same skew & 2 & 0.0000 & 0.0000 & 0.0000 & 0.0000 \\
2x2 & same skew & 5 & 0.1110 & 0.7282 & 0.0000 & 0.2304 \\
2x2 & same skew & 10 & 0.0357 & 0.9585 & 0.0092 & 0.6511 \\
2x2 & same skew & 20 & 0.0232 & 0.9980 & 0.0135 & 0.9550 \\
2x2 & same skew & 50 & 0.0086 & 1.0000 & 0.0037 & 1.0000 \\
2x2 & same skew & 100 & 0.0107 & 1.0000 & 0.0047 & 1.0000 \\
2x2 & same skew & 250 & 0.0260 & 1.0000 & 0.0135 & 1.0000 \\
2x2 & same skew & 500 & 0.0379 & 1.0000 & 0.0268 & 1.0000 \\
2x2 & same skew & 1000 & 0.0437 & 1.0000 & 0.0372 & 1.0000 \\
2x2 & uniform & 2 & 0.0000 & 0.0000 & 0.0000 & 0.0000 \\
2x2 & uniform & 5 & 0.0914 & 0.9862 & 0.0000 & 0.7774 \\
2x2 & uniform & 10 & 0.0588 & 0.9889 & 0.0263 & 0.9795 \\
2x2 & uniform & 20 & 0.0261 & 0.9977 & 0.0164 & 0.9833 \\
2x2 & uniform & 50 & 0.0082 & 0.9999 & 0.0045 & 0.9938 \\
2x2 & uniform & 100 & 0.0142 & 1.0000 & 0.0063 & 0.9991 \\
2x2 & uniform & 250 & 0.0306 & 1.0000 & 0.0164 & 1.0000 \\
2x2 & uniform & 500 & 0.0408 & 1.0000 & 0.0298 & 1.0000 \\
2x2 & uniform & 1000 & 0.0445 & 1.0000 & 0.0386 & 1.0000 \\
3x3 & different skew & 2 & 0.0000 & 0.0000 & 0.0000 & 0.0000 \\
3x3 & different skew & 5 & 0.2089 & 0.8383 & 0.0000 & 0.3296 \\
3x3 & different skew & 10 & 0.0886 & 0.9886 & 0.0260 & 0.7792 \\
3x3 & different skew & 20 & 0.0466 & 0.9998 & 0.0251 & 0.9745 \\
3x3 & different skew & 50 & 0.0188 & 1.0000 & 0.0098 & 1.0000 \\
3x3 & different skew & 100 & 0.0171 & 1.0000 & 0.0083 & 1.0000 \\
3x3 & different skew & 250 & 0.0268 & 1.0000 & 0.0147 & 1.0000 \\
3x3 & different skew & 500 & 0.0331 & 1.0000 & 0.0232 & 1.0000 \\
3x3 & different skew & 1000 & 0.0411 & 1.0000 & 0.0348 & 1.0000 \\
3x3 & same skew & 2 & 0.0000 & 0.0000 & 0.0000 & 0.0000 \\
3x3 & same skew & 5 & 0.1807 & 0.7217 & 0.0000 & 0.2209 \\
3x3 & same skew & 10 & 0.0602 & 0.9599 & 0.0176 & 0.6627 \\
3x3 & same skew & 20 & 0.0333 & 0.9996 & 0.0207 & 0.9565 \\
3x3 & same skew & 50 & 0.0135 & 1.0000 & 0.0076 & 1.0000 \\
3x3 & same skew & 100 & 0.0139 & 1.0000 & 0.0068 & 1.0000 \\
3x3 & same skew & 250 & 0.0273 & 1.0000 & 0.0149 & 1.0000 \\
3x3 & same skew & 500 & 0.0348 & 1.0000 & 0.0235 & 1.0000 \\
3x3 & same skew & 1000 & 0.0391 & 1.0000 & 0.0314 & 1.0000 \\
3x3 & uniform & 2 & 0.0000 & 0.0000 & 0.0000 & 0.0000 \\
3x3 & uniform & 5 & 0.1931 & 0.9994 & 0.0000 & 0.9534 \\
3x3 & uniform & 10 & 0.1058 & 1.0000 & 0.0285 & 0.9981 \\
3x3 & uniform & 20 & 0.0730 & 1.0000 & 0.0343 & 1.0000 \\
3x3 & uniform & 50 & 0.0288 & 1.0000 & 0.0173 & 1.0000 \\
3x3 & uniform & 100 & 0.0167 & 1.0000 & 0.0106 & 1.0000 \\
3x3 & uniform & 250 & 0.0269 & 1.0000 & 0.0178 & 1.0000 \\
3x3 & uniform & 500 & 0.0355 & 1.0000 & 0.0278 & 1.0000 \\
3x3 & uniform & 1000 & 0.0445 & 1.0000 & 0.0393 & 1.0000 \\
5x5 & different skew & 2 & 0.0000 & 0.0000 & 0.0000 & 0.0000 \\
5x5 & different skew & 5 & 0.2690 & 0.8356 & 0.0000 & 0.3270 \\
5x5 & different skew & 10 & 0.1514 & 0.9889 & 0.0693 & 0.7817 \\
5x5 & different skew & 20 & 0.1259 & 1.0000 & 0.0839 & 0.9785 \\
5x5 & different skew & 50 & 0.0623 & 1.0000 & 0.0348 & 1.0000 \\
5x5 & different skew & 100 & 0.0250 & 1.0000 & 0.0136 & 1.0000 \\
5x5 & different skew & 250 & 0.0254 & 1.0000 & 0.0162 & 1.0000 \\
5x5 & different skew & 500 & 0.0360 & 1.0000 & 0.0268 & 1.0000 \\
5x5 & different skew & 1000 & 0.0386 & 1.0000 & 0.0319 & 1.0000 \\
5x5 & same skew & 2 & 0.0000 & 0.0000 & 0.0000 & 0.0000 \\
5x5 & same skew & 5 & 0.2175 & 0.7137 & 0.0000 & 0.2149 \\
5x5 & same skew & 10 & 0.0825 & 0.9599 & 0.0363 & 0.6756 \\
5x5 & same skew & 20 & 0.0651 & 0.9996 & 0.0454 & 0.9553 \\
5x5 & same skew & 50 & 0.0290 & 1.0000 & 0.0178 & 1.0000 \\
5x5 & same skew & 100 & 0.0112 & 1.0000 & 0.0066 & 1.0000 \\
5x5 & same skew & 250 & 0.0152 & 1.0000 & 0.0094 & 1.0000 \\
5x5 & same skew & 500 & 0.0369 & 1.0000 & 0.0260 & 1.0000 \\
5x5 & same skew & 1000 & 0.0403 & 1.0000 & 0.0329 & 1.0000 \\
5x5 & uniform & 2 & 0.0000 & 0.0000 & 0.0000 & 0.0000 \\
5x5 & uniform & 5 & 0.2287 & 1.0000 & 0.0000 & 0.9909 \\
5x5 & uniform & 10 & 0.1061 & 1.0000 & 0.0665 & 0.9997 \\
5x5 & uniform & 20 & 0.0768 & 1.0000 & 0.0528 & 1.0000 \\
5x5 & uniform & 50 & 0.0643 & 1.0000 & 0.0467 & 1.0000 \\
5x5 & uniform & 100 & 0.0377 & 1.0000 & 0.0289 & 1.0000 \\
5x5 & uniform & 250 & 0.0232 & 1.0000 & 0.0190 & 1.0000 \\
5x5 & uniform & 500 & 0.0276 & 1.0000 & 0.0240 & 1.0000 \\
5x5 & uniform & 1000 & 0.0326 & 1.0000 & 0.0296 & 1.0000 \\
8x8 & different skew & 2 & 0.0000 & 0.0000 & 0.0000 & 0.0000 \\
8x8 & different skew & 5 & 0.2986 & 0.8319 & 0.0000 & 0.3157 \\
8x8 & different skew & 10 & 0.1853 & 0.9897 & 0.1035 & 0.7912 \\
8x8 & different skew & 20 & 0.1803 & 1.0000 & 0.1432 & 0.9759 \\
8x8 & different skew & 50 & 0.1391 & 1.0000 & 0.1080 & 1.0000 \\
8x8 & different skew & 100 & 0.0921 & 1.0000 & 0.0699 & 1.0000 \\
8x8 & different skew & 250 & 0.0583 & 1.0000 & 0.0449 & 1.0000 \\
8x8 & different skew & 500 & 0.0377 & 1.0000 & 0.0300 & 1.0000 \\
8x8 & different skew & 1000 & 0.0356 & 1.0000 & 0.0301 & 1.0000 \\
8x8 & same skew & 2 & 0.0000 & 0.0000 & 0.0000 & 0.0000 \\
8x8 & same skew & 5 & 0.2266 & 0.6984 & 0.0000 & 0.2082 \\
8x8 & same skew & 10 & 0.0873 & 0.9587 & 0.0456 & 0.6796 \\
8x8 & same skew & 20 & 0.0873 & 0.9992 & 0.0692 & 0.9541 \\
8x8 & same skew & 50 & 0.0537 & 1.0000 & 0.0367 & 1.0000 \\
8x8 & same skew & 100 & 0.0245 & 1.0000 & 0.0155 & 1.0000 \\
8x8 & same skew & 250 & 0.0106 & 1.0000 & 0.0071 & 1.0000 \\
8x8 & same skew & 500 & 0.0162 & 1.0000 & 0.0129 & 1.0000 \\
8x8 & same skew & 1000 & 0.0341 & 1.0000 & 0.0289 & 1.0000 \\
8x8 & uniform & 2 & 0.0000 & 0.0000 & 0.0000 & 0.0000 \\
8x8 & uniform & 5 & 0.2692 & 0.9983 & 0.0000 & 0.9163 \\
8x8 & uniform & 10 & 0.1376 & 1.0000 & 0.0762 & 0.9996 \\
8x8 & uniform & 20 & 0.0687 & 1.0000 & 0.0570 & 1.0000 \\
8x8 & uniform & 50 & 0.0703 & 1.0000 & 0.0592 & 1.0000 \\
8x8 & uniform & 100 & 0.0712 & 1.0000 & 0.0618 & 1.0000 \\
8x8 & uniform & 250 & 0.0317 & 1.0000 & 0.0286 & 1.0000 \\
8x8 & uniform & 500 & 0.0212 & 1.0000 & 0.0198 & 1.0000 \\
8x8 & uniform & 1000 & 0.0228 & 1.0000 & 0.0211 & 1.0000 \\
\bottomrule\end{longtable}
\endgroup

